\section{Введение}

\textbf{Актуальность.} Генерация с дополненным поиском (Retrieval-Augmented Generation, RAG\footnote{RAG --- схема, в которой ответ модели строится по фрагментам документов, найденным поиском по запросу пользователя~\cite{lewis2020}.}) стала стандартной архитектурой вопросно-ответных систем над корпоративными и открытыми базами знаний. Качество такой системы зависит от набора параметров: способа поиска, числа фрагментов, передаваемых модели (top-k), размера фрагмента (чанка) и числа шагов поиска. На практике эти параметры подбирают по средним значениям метрик на небольших наборах вопросов, без проверки того, что разница между конфигурациями не случайна. Сравнения систем поиска без статистических тестов --- известная проблема в информационном поиске и NLP~\cite{smucker2007,dror2018}.

\textbf{Цель} --- статистически оценить, какие параметры RAG и какие свойства вопроса значимо влияют на качество поиска и ответа, и насколько правильность ответа предсказуема по логам.

\textbf{Задачи:}
\begin{enumerate}
\item выбрать и обосновать датасеты из пяти заданных (MS MARCO, Natural Questions, HotpotQA, BEIR, SQuAD);
\item прогнать систему ретривер\,+\,ридер по сетке параметров и записать логи поиска и ответов;
\item получить описательную статистику метрик по конфигурациям;
\item проверить гипотезы о влиянии параметров и свойств вопроса на качество;
\item оценить корреляции параметров и метрик;
\item построить регрессию качества по параметрам и классификатор правильности ответа;
\item оценить динамику качества по итерациям поиска.
\end{enumerate}

\textbf{Гипотезы.}
\begin{itemize}
\item[H1.] Тип ретривера (BM25, плотный, гибридный) значимо влияет на полноту поиска и на F1 ответа.
\item[H2.] С ростом top-k полнота поиска монотонно растет, а F1 ответа растет медленнее и выходит на плато.
\item[H3.] Размер чанка значимо влияет на качество ответа.
\item[H4.] На многошаговых вопросах качество ниже на вопросах типа bridge, чем comparison, и падает с ростом уровня сложности.
\item[H5.] Наличие ответа в найденном контексте повышает шанс правильного ответа, но не гарантирует его.
\item[H6.] Итеративный поиск повышает полноту на многошаговых вопросах по сравнению с одношаговым поиском того же объема контекста и не дает выигрыша на одношаговых.
\end{itemize}

\section{Обзор датасетов}

Преподавателем заданы пять наборов данных; дополнительно рассмотрены три близких набора для многошаговых и открытых вопросов. Характеристики --- в таблице~\ref{tab:datasets}. BEIR включает в себя Natural Questions, HotpotQA и MS MARCO как подзадачи поиска, но без ответов, поэтому он рассматривается как источник разметки релевантности, а не как отдельный корпус.

\begin{gosttable}{Рассмотренные датасеты: пять заданных и три дополнительных}\label{tab:datasets}
\begin{tabular}{>{\raggedright\arraybackslash}p{2.5cm}>{\raggedright\arraybackslash}p{3.4cm}>{\raggedright\arraybackslash}p{4.1cm}>{\raggedright\arraybackslash}p{4.1cm}}
\toprule
Датасет & Объем & Что размечено & Ограничение для задачи \\
\midrule
MS MARCO~\cite{msmarco} & 1\,010\,916 запросов Bing, 8,8~млн пассажей & 10 пассажей выдачи Bing на запрос с отметкой полезности, свободный ответ человека & ответы перефразированы, F1 к ним занижен у экстрактивного ридера; корпус тяжелый \\
Natural Questions~\cite{nq} & 307\,373 обучающих вопроса, 7\,830 dev & реальные запросы Google, длинный и короткий ответ на странице Википедии & корпус --- полные страницы, архив около 40~ГБ \\
HotpotQA~\cite{hotpotqa} & 90\,447 train, 7\,405 dev (distractor) & 2 нужных абзаца и 8 отвлекающих, опорные предложения, тип bridge/comparison, уровень easy/medium/hard & в dev все вопросы уровня hard \\
BEIR~\cite{beir} & 18 наборов, 9 типов задач & запросы и релевантность документов & нет ответов --- только поиск \\
SQuAD v1.1~\cite{squad} & 87\,599 train, 10\,570 dev, 536 статей & ответ --- отрезок абзаца Википедии & вопросы написаны по абзацу, высокое лексическое пересечение \\
\midrule
2Wiki\-Multihop\-QA~\cite{twowiki} & 192\,606 вопросов & 2--4 нужных абзаца, цепочка фактов из Викиданных & вопросы шаблонные, вне заданного списка \\
MuSiQue~\cite{musique} & около 25 тыс. вопросов на 2--4 шага & нужные абзацы, разбиение на подвопросы & вне заданного списка, вопросы собраны из одношаговых \\
TriviaQA~\cite{triviaqa} & 95 тыс. пар вопрос--ответ, 650 тыс. документов & ответ и найденные документы без разметки нужного абзаца & нет разметки релевантности --- не измерить поиск \\
\bottomrule
\end{tabular}
\end{gosttable}

Для анализа выбраны \textbf{HotpotQA} и \textbf{SQuAD v1.1} из заданного списка.
\begin{itemize}
\item В обоих есть и разметка релевантности (нужные абзацы и опорные предложения у HotpotQA, исходный абзац у SQuAD), и эталонный ответ. Это позволяет отдельно мерить поиск и ответ в одном логе.
\item Ответы в обоих экстрактивные --- это отрезок текста, поэтому EM и F1 корректны без LLM-судьи.
\item HotpotQA дает факторы вопроса для гипотез: тип (bridge --- ответ через промежуточную сущность, comparison --- сравнение двух объектов) и уровень сложности.
\item SQuAD --- одношаговые вопросы. Он служит контрольной группой для итеративного поиска и для эффекта лексического пересечения.
\item Оба набора умещаются в память ноутбука: прогон воспроизводим без GPU-сервера.
\end{itemize}

MS MARCO и Natural Questions не выбраны из-за объема корпуса и свободной формы ответов. BEIR не содержит ответов, поэтому на нем не измерить вторую половину RAG.

\section{Обзор методов}

\textbf{Схема ретривер--ридер.} Chen и соавторы~\cite{drqa} предложили для открытых вопросов по Википедии двухступенчатую схему: TF-IDF-поиск абзацев и нейросетевой ридер, выделяющий ответ. Lewis и соавторы~\cite{lewis2020} заменили ридер генератором и ввели название RAG. Обзор~\cite{gao2023} выделяет параметры, которые подбирают на практике: тип ретривера, размер чанка, top-k, итеративный поиск. Эта работа использует ту же схему с экстрактивным ридером.

\textbf{Ретриверы.} BM25~\cite{bm25} --- лексическая модель ранжирования, базовая линия во всех бенчмарках поиска. Плотный поиск DPR~\cite{dpr} на Natural Questions превзошел BM25 по точности top-20 на 9--19 процентных пунктов. Эмбеддинги E5~\cite{e5} обучены контрастивно на слабой разметке и сильны без дообучения. Слияние выдач через RRF~\cite{rrf}\footnote{Reciprocal Rank Fusion: документ получает сумму $1/(60+r_i)$ по рангам $r_i$ в каждой выдаче.} устойчиво превосходит отдельные системы. BEIR~\cite{beir} показал, что вне домена обучения BM25 часто сильнее плотных моделей. Поэтому тип ретривера --- фактор с неочевидным знаком эффекта.

\textbf{Оценка RAG.} RAGAS~\cite{ragas} и ARES~\cite{ares} оценивают верность ответа контексту и релевантность контекста LLM-судьей, ARES --- с доверительными интервалами по малой размеченной выборке. eRAG~\cite{erag} показал, что классические метрики поиска по разметке релевантности слабо коррелируют с итоговым качеством RAG, и предложил оценивать документ по тому, помогает ли он ответить. Здесь та же связь проверяется прямой корреляцией полноты поиска и F1 ответа.

\textbf{Влияние контекста.} Liu и соавторы~\cite{lostmiddle} показали, что качество ответа падает, когда нужный фрагмент в середине длинного контекста, и что рост числа документов перестает помогать. Cuconasu и соавторы~\cite{powernoise} измерили влияние отвлекающих документов на точность RAG на Natural Questions. Это основа гипотезы H2 о плато.

\textbf{Итеративный поиск.} IRCoT~\cite{ircot} чередует шаги поиска и рассуждения и на четырех многошаговых наборах, включая HotpotQA, повышает полноту поиска до 21 пункта против одношагового поиска. В этой работе используется упрощенный вариант без LLM: к запросу дописывается промежуточный ответ ридера.

\textbf{Статистическое сравнение систем.} Smucker и соавторы~\cite{smucker2007} сравнили тесты значимости для оценки поиска и рекомендовали парные тесты на одних и тех же запросах. Dror и соавторы~\cite{dror2018} дали схему выбора теста для метрик NLP: для ненормальных и ограниченных метрик --- непараметрические парные тесты, при множественных сравнениях --- поправка. Отсюда методика этой работы: тест Фридмана~\cite{friedman} и Вилкоксона для F1 и полноты, Q Кохрена~\cite{cochran} и Мак-Немара~\cite{mcnemar} для бинарной правильности, тренд Пейджа~\cite{page} для упорядоченных уровней, поправка Холма~\cite{holm}, бутстреп-интервалы~\cite{efron}. Наблюдения одного вопроса в разных конфигурациях зависимы, поэтому регрессия считается с кластерными ошибками, а логистическая --- методом обобщенных оценочных уравнений (GEE~\cite{gee}).

\section{Методология и пайплайн анализа}

\subsection{Данные и выборка}

Из HotpotQA (distractor, train: 17\,972 вопроса easy, 56\,814 medium, 15\,661 hard) взято по 200 вопросов каждого уровня без ответов yes/no\footnote{Ответы yes/no не встречаются в тексте как отрезок, экстрактивный ридер их не выдает; таких вопросов в train 6\,\%.}. Корпус --- все абзацы выбранных вопросов: два нужных и восемь отвлекающих на вопрос, после удаления повторов по заголовку. Отвлекающие абзацы в HotpotQA подобраны поиском по биграммам, то есть корпус заведомо содержит лексически похожие документы. Из SQuAD v1.1 dev взято 600 случайных вопросов, корпус --- все 2\,067 абзацев dev-части. Зерно выборки --- 42.

\subsection{Система RAG и сетка параметров}

\begin{itemize}
\item \textbf{Чанки:} документы режутся по предложениям до 50, 100 или 200 слов, к чанку приписывается заголовок документа.
\item \textbf{Ретриверы:} BM25 (bm25s, $k_1=1{,}5$, $b=0{,}75$), плотный intfloat/e5-small-v2, гибрид через RRF по выдачам глубины 50.
\item \textbf{top-k:} 1, 3, 5, 10.
\item \textbf{Ридер:} deepset/roberta-base-squad2. На каждый чанк выдачи ищется лучший отрезок, ответ системы --- отрезок с наибольшей оценкой среди top-k чанков.
\end{itemize}

Сетка: 2 датасета $\times$ 3 ретривера $\times$ 3 размера чанка $\times$ 4 значения top-k = 72 конфигурации, по 600 вопросов в каждой --- 43\,200 строк лога ответов. Лог поиска хранит выдачу глубины 10 с оценкой ретривера, отметкой релевантности и ответом ридера на каждый чанк.

\textbf{Итеративный поиск} (гибрид, чанк 100): на шаге $t$ в контекст добавляются 3 новых чанка, к запросу дописывается ответ ридера по новым чанкам. Контроль --- одношаговый поиск того же объема (top-$3t$). Шагов 4.

\subsection{Метрики}

\begin{itemize}
\item \textbf{Релевантный чанк}: у HotpotQA --- чанк нужного абзаца, содержащий опорное предложение; у SQuAD --- чанк исходного абзаца, содержащий ответ (если ответ разрезан границей чанков --- все чанки абзаца).
\item \textbf{Полнота поиска} (recall@k) --- доля нужных документов вопроса, релевантный чанк которых попал в top-k (у HotpotQA нужных два, у SQuAD один).
\item \textbf{Попадание} (hit@k) --- в top-k есть хотя бы один релевантный чанк. \textbf{MRR@k} --- обратный ранг первого релевантного чанка.
\item \textbf{Ответ в контексте} --- эталонный ответ встречается подстрокой в top-k чанках.
\item \textbf{EM} --- точное совпадение после нормализации SQuAD. \textbf{F1} --- гармоническое среднее точности и полноты по словам ответа~\cite{squad}.
\item \textbf{Правильный ответ} --- F1~$\ge 0{,}5$.
\item \textbf{Лексическое пересечение} --- доля значимых слов вопроса (без стоп-слов), встречающихся в эталонных абзацах; считается по эталону и описывает вопрос, а не систему.
\end{itemize}

\subsection{Пайплайн анализа}

\begin{enumerate}
\item Описательная статистика по конфигурациям, бутстреп-интервалы (2\,000 повторов).
\item Парные тесты по вопросам: Фридман, Вилкоксон (нулевые разности отбрасываются) для F1, Q Кохрена и Мак-Немар для бинарных метрик, тренд Пейджа, поправка Холма, бутстреп-интервал парной разности и TOST для проверки плато. Межгрупповые: Манн--Уитни, Краскел--Уоллис, $\chi^2$.
\item Корреляции Спирмена по всей сетке, проверка значимости ключевых пар.
\item OLS для F1 с кластерными ошибками по вопросу, GEE-логит для правильности, VIF и остатки.
\item Классификация правильности по признакам без эталона: логистическая регрессия и градиентный бустинг, 5-кратная групповая кросс-валидация по вопросу.
\item Динамика по итерациям: Фридман и тренд Пейджа, Вилкоксон итеративный против одношагового, GEE с взаимодействием итерация $\times$ режим.
\end{enumerate}

Структура логов --- в приложении~\ref{app:logs}, команды воспроизведения --- в приложении~\ref{app:repro}.
